%-----------------------------------------------------------------------------%
\chapter{\babDua}
\label{cha:studiliteratur}
%-----------------------------------------------------------------------------%
Bab ini menyajikan teori dan penelitian terkait yang diperlukan untuk merancang serta mengevaluasi DroneVL. Pembahasan dimulai dari navigasi UAV dan kemampuan model VLM dan VLA, kemudian menjelaskan penambatan visual, penalaran spasial, navigasi \textit{closed-loop}, integrasi CoSyS-AirSim--ROS~2, kontrak aksi, dan evaluasi yang dapat direproduksi. Bagian akhir membandingkan pendekatan terdahulu dengan kontribusi DroneVL untuk mengidentifikasi celah penelitian yang diuji.

%-----------------------------------------------------------------------------%
\section{Landasan Teoretis}
\label{sec:landasan_teori}
%-----------------------------------------------------------------------------%

\subsection{UAV dan Navigasi Otonom}
\label{subsec:uav_navigasi}
UAV adalah kendaraan udara yang dapat beroperasi tanpa pilot di dalam kendaraan. Pada multirotor, gaya dorong beberapa rotor menghasilkan gerak translasi dan rotasi, sedangkan tumpukan navigasi mengestimasi keadaan kendaraan, membentuk sasaran, dan meneruskan target posisi, kecepatan, ketinggian, atau yaw kepada pengendali penerbangan. Keputusan selama penerbangan perlu mempertimbangkan informasi yang diperoleh dari lingkungan karena ketidakpastian dapat mengubah lintasan yang layak ditempuh \citep{vivaldini2025decision}. Pada navigasi tujuan-objek, kebutuhan tersebut mencakup pencarian, pengenalan, dan pelokalan objek sasaran sebelum kendaraan dapat membentuk rute menuju sasaran \citep{ayala2024uav}. % [SRC-B2-01]

Navigasi geometris umumnya menyatakan tujuan melalui koordinat, \textit{waypoint}, peta, atau representasi spasial lain yang dapat langsung dipakai perencana. Navigasi semantik menyatakan tujuan sebagai objek, atribut, relasi, atau instruksi; sistem harus terlebih dahulu menambatkan deskripsi tersebut, kemudian membentuk satu atau lebih subtujuan geometris. Kedua pendekatan tidak saling meniadakan: DroneVL menggunakan penalaran semantik untuk menentukan sasaran atau relasi yang relevan, lalu menggunakan kendali dan perencanaan geometris untuk mencapai subtujuan secara aman. Perbedaan ini penting karena objek dapat tertutup, berada di luar sudut pandang awal, atau memiliki beberapa kandidat yang serupa. % [SRC-B2-02]

\subsection{\textit{Vision-Language Model}}
\label{subsec:vlm}
\textit{Vision-language model} (VLM) menghubungkan representasi visual dan bahasa melalui penyandi visual, model bahasa, serta mekanisme penyelarasan lintasmodal. CLIP memperlihatkan bahwa supervisi bahasa alami dapat menghasilkan representasi gambar--teks yang dapat ditransfer ke tugas visual, sehingga pendekatan tersebut tidak terbatas pada daftar kelas objek yang tetap \citep{radford2021clip}. Penelitian \textit{Show and Tell} juga menunjukkan bahwa arsitektur saraf dapat mengaitkan visi komputer dengan pemrosesan bahasa alami untuk menghasilkan deskripsi gambar \citep{vinyals2015show}. Landasan ini mendukung penggunaan VLM sebagai komponen penalaran semantik, tetapi tidak menjadikan keluaran bahasa sebagai bukti fisik atau perintah penerbangan yang aman. % [SRC-B2-03]

Keluaran VLM bersifat probabilistik dan peka terhadap resolusi, pemotongan gambar, urutan konteks, templat \textit{prompt}, dan parameter dekode. Model dapat salah mengenali objek, menghalusinasi objek yang tidak terlihat, atau menghasilkan respons yang tidak sesuai dengan tindakan yang tersedia. Oleh karena itu, DroneVL memperlakukan VLM sebagai komponen inferensi dengan kontrak masukan dan keluaran yang ketat. Kontrak tersebut membatasi perbedaan antarmodel pada adaptor, sementara simulator, pengendali, evaluator, dan kebijakan keselamatan tetap menggunakan representasi bersama. Dengan pemisahan ini, kesalahan respons dapat ditelusuri tanpa menyamakan respons bahasa model dengan keadaan sebenarnya dari lingkungan.

\subsection{\textit{Vision-Language-Action}}
\label{subsec:vla}
Model \textit{vision-language-action} (VLA) menghubungkan pengamatan visual dan instruksi bahasa dengan keluaran aksi kebijakan. OpenVLA-7B adalah VLA terbuka yang menghasilkan token aksi diskret dan menyediakan prosedur \textit{fine-tuning} LoRA untuk tugas serta perwujudan robot baru \citep{kim2024openvla}. OpenPI menyediakan keluarga $\pi_0$, $\pi_0$-FAST, dan $\pi_{0.5}$; implementasi $\pi_{0.5}$ saat ini menggunakan kepala \textit{flow-matching} untuk menghasilkan potongan aksi kontinu \citep{physicalintelligence2025openpi}. CognitiveDrone berbeda karena dirancang khusus untuk UAV, dilatih pada trajektori penerbangan simulasi, dan menghasilkan perintah aksi empat-dimensi dari citra orang-pertama serta instruksi \citep{lykov2025cognitivedrone}.

Keluaran aksi VLA selalu terikat pada perwujudan dan data pelatihan asalnya. Karena OpenVLA-7B dan OpenPI $\pi_{0.5}$ ditujukan untuk manipulasi robot, aksi sendi atau potongan aksi asal tidak dapat dipetakan langsung menjadi kendali UAV. DroneVL hanya mengevaluasi versi yang telah diadaptasi terhadap skema aksi UAV CoSyS-AirSim melalui trajektori berpasangan RGB--instruksi--keadaan UAV--aksi UAV. Setiap skema menentukan satuan, kerangka koordinat, horizon aksi, dan prosedur dekode sebelum keluaran dapat menjadi kandidat aksi kanonis (\textit{canonical action}). Batas ini membedakan adaptasi kebijakan yang dapat diuji dari klaim transfer lintas-perwujudan \textit{zero-shot}.

\subsection{Pemilihan Model untuk Eksperimen}
\label{subsec:pemilihan_model}
Pemilihan model pada DroneVL bersifat bertujuan, bukan peringkat kemampuan umum. Model dipilih apabila dapat menerima pengamatan visual dan instruksi, menyediakan antarmuka keluaran yang dapat dinormalisasi ke kontrak DroneVL, serta mewakili variasi yang relevan bagi pertanyaan penelitian: akses layanan tertutup atau model yang dapat dikonfigurasi, keluaran terstruktur atau pemanggilan fungsi, penalaran berwujud, dan kebijakan aksi yang terikat pada perwujudan. Kapabilitas yang dideklarasikan penyedia dicatat sebagai konfigurasi awal; keunggulan pada penambatan visual, keselamatan, efisiensi lintasan, dan latensi hanya disimpulkan dari eksperimen pada Bab~III dan Bab~IV.

\subsubsection{Kelompok VLM}
GPT-5.6 dipilih sebagai pembanding VLM berbasis API untuk penalaran multimodal tingkat tinggi. Dokumentasi OpenAI menyatakan bahwa alias \texttt{gpt-5.6} merujuk ke GPT-5.6 Sol dan mendukung masukan citra, keluaran terstruktur, serta pemanggilan fungsi \citep{openai2026gpt56}. Karakteristik ini memungkinkan keluaran perencana dikirim dalam format yang dapat diperiksa oleh pengurai DroneVL. Model ini tidak digunakan sebagai pengendali penerbangan, melainkan sebagai sumber keputusan semantik yang seluruh keluarannya tetap melewati normalisasi dan gerbang keselamatan.

Qwen3.6 35B-A3B dipilih sebagai VLM \textit{mixture-of-experts} (MoE) multimodal yang konfigurasi \textit{runtime} dan prapemrosesannya dapat didokumentasikan secara eksplisit pada manifest eksperimen. Varian yang digunakan adalah \texttt{qwen3.6:35b-a3b} melalui Ollama; halaman model mencatat bobot 35,5B, arsitektur \texttt{qwen35moe}, kuantisasi Q4\_K\_M, dan proyektor CLIP \citep{ollama2026qwen36}. Dokumentasi Qwen juga mencantumkan dukungan masukan teks, citra, dan video, pemanggilan fungsi, serta keluaran terstruktur untuk varian 35B-A3B \citep{qwen2026visionmodels}. Karena itu, Qwen3.6 35B-A3B memungkinkan pengujian apakah kontrak pengamatan dan aksi DroneVL tetap berlaku ketika model, cara penyajian citra, dan lingkungan inferensi berbeda dari layanan API tertutup. Tag model, checksum, kuantisasi, dan perangkat keras yang benar-benar digunakan dikunci sebelum pengujian.

Gemini Robotics ER 2 dengan varian \textit{Streaming Preview} dipilih karena dirancang untuk penalaran berwujud, termasuk penalaran spasial, pemahaman kemajuan video, dan orkestrasi tugas multilangkah. Varian ini juga dioptimalkan untuk interaksi dua arah berlatensi rendah melalui Live API dan mendukung pemanggilan fungsi \citep{google2026geminirobotics}. Dengan demikian, Gemini menguji apakah adaptor DroneVL dapat menangani respons \textit{streaming} dari VLM yang berorientasi robotika tanpa mengubah pengendali UAV. Fokus evaluasinya tetap pada kualitas keluaran yang berhasil dinormalisasi dan dieksekusi secara aman dalam simulator, bukan pada klaim kemampuan robotika penyedia.

Ketiga VLM tersebut membentuk pembandingan antarmuka yang saling melengkapi: GPT-5.6 sebagai layanan API multimodal dengan keluaran terstruktur, Qwen3.6 35B-A3B sebagai VLM MoE yang dijalankan secara lokal dengan konfigurasi inferensi yang dapat dikendalikan dan dicatat, serta Gemini sebagai VLM berorientasi robotika dengan jalur \textit{streaming}. Perbedaan akses, modalitas, dan antarmuka sengaja tidak dihapus; seluruhnya dilaporkan bersama hasil agar perbandingan tidak menganggap ketiga model memiliki kondisi operasi yang identik.

\subsubsection{Kelompok VLA}
CognitiveDrone dipilih sebagai titik acuan VLA yang sejak awal ditujukan bagi UAV. Model tersebut menggunakan pengamatan orang-pertama dan instruksi untuk menghasilkan aksi empat-dimensi pada trajektori penerbangan simulasi \citep{lykov2025cognitivedrone}. Perannya adalah menguji adaptor pada kebijakan yang sudah dekat dengan perwujudan UAV, sambil tetap mewajibkan deklarasi skema aksi, satuan, kerangka koordinat, dan horizon keluaran.

OpenVLA-7B dipilih sebagai VLA utama untuk eksperimen adaptasi karena keluaran aksinya berupa token diskret dan implementasinya menyediakan prosedur \textit{fine-tuning} LoRA untuk tugas serta perwujudan robot baru \citep{kim2024openvla}. Sifat tersebut mendukung tiga kondisi penelitian: model dasar dengan tokenisasi aksi UAV (M0), adaptasi LoRA pada trajektori CoSyS-AirSim (M1), dan M1 yang dilanjutkan GRPO (M2). Dengan kata lain, OpenVLA-7B bukan hanya pembanding portabilitas, tetapi objek utama untuk menguji pengaruh LoRA dan GRPO pada kebijakan UAV yang menggunakan kontrak aksi yang sama.

OpenPI $\pi_{0.5}$ dipilih untuk menguji batas portabilitas adaptor terhadap VLA dengan representasi aksi yang berbeda. OpenPI menyediakan kebijakan yang menghasilkan potongan aksi kontinu melalui kepala \textit{flow-matching} \citep{physicalintelligence2025openpi}; setelah adaptasi UAV, keluarannya harus diterjemahkan melalui \textit{NativeActionEnvelope} sebelum menjadi kandidat aksi kanonis. Model ini tidak dimasukkan ke eksperimen GRPO, karena keluaran kontinu tersebut tidak serta-merta menyediakan rasio probabilitas kebijakan yang diperlukan oleh prosedur GRPO. Perannya adalah menguji kontrak dan pengaman DroneVL pada perbedaan representasi aksi, bukan membandingkan metode optimisasi yang tidak sepadan.

Secara bersama, kelompok VLA mencakup kebijakan yang dirancang untuk UAV, VLA bertoken diskret yang menjadi objek adaptasi utama, dan VLA dengan aksi kontinu. Susunan ini memungkinkan penelitian membedakan dua pertanyaan: apakah adaptor dapat memportasikan model dengan asumsi perwujudan serta format aksi yang berbeda, dan apakah LoRA+GRPO pada OpenVLA-7B yang telah diadaptasi meningkatkan kinerja terhadap kondisi dasarnya. Kedua pertanyaan dilaporkan terpisah agar hasil portabilitas tidak ditafsirkan sebagai efek LoRA atau GRPO.

\subsection{\textit{Low-Rank Adaptation} dan \textit{Group Relative Policy Optimization}}
\label{subsec:lora_grpo}
\textit{Low-rank adaptation} (LoRA) adalah pendekatan \textit{parameter-efficient fine-tuning} yang membekukan bobot model pralatih dan menambahkan matriks berperingkat rendah yang dapat dilatih pada lapisan transformator \citep{hu2021lora}. Pendekatan tersebut relevan bagi model terbuka karena adaptasi domain tidak memerlukan pembaruan seluruh parameter. QLoRA menggabungkan LoRA dengan kuantisasi empat-bit pada model pralatih yang dibekukan untuk mengurangi kebutuhan memori selama \textit{fine-tuning} \citep{dettmers2023qlora}. Dalam DroneVL, LoRA dan QLoRA hanya diterapkan pada OpenVLA-7B sebagai VLA utama; pendekatan ini tidak diterapkan pada VLM API tertutup maupun pada Qwen3.6 35B-A3B agar efek adaptasi kebijakan aksi UAV dapat diisolasi. % [SRC-B2-10]

\textit{Group Relative Policy Optimization} (GRPO) diperkenalkan sebagai varian \textit{proximal policy optimization} yang menggunakan perbandingan hasil dalam satu kelompok respons dan ditujukan untuk menurunkan kebutuhan memori dibandingkan pendekatan dengan \textit{critic} terpisah \citep{shao2024deepseekmath}. Dalam ranah multimodal, Vision-R1 menerapkan GRPO bersama penghargaan format dan data penalaran multimodal, sedangkan VLM-R1 mengeksplorasi pembelajaran penguatan bergaya R1 pada tugas visual dengan \textit{ground truth} yang terdefinisi \citep{huang2025visionr1,shen2025vlmr1}. Pada DroneVL, eksperimen utama LoRA+GRPO diterapkan pada OpenVLA-7B yang telah diberi tokenisasi aksi UAV; keluaran aksi diskretnya memungkinkan probabilitas token aksi dari setiap \textit{rollout} dicatat untuk rasio kebijakan GRPO. Tiga kondisi dibandingkan pada set pengujian yang sama dan telah dibekukan: model dasar OpenVLA-7B (M0), OpenVLA-7B dengan LoRA yang diadaptasi untuk UAV (M1), dan M1 yang dilanjutkan GRPO pada \textit{rollout} CoSyS-AirSim (M2). Fungsi penghargaan, data pelatihan, anggaran \textit{rollout}, komputasi, dan komponen ablasinya ditetapkan di Bab~III sebelum pelatihan dimulai. OpenPI $\pi_{0.5}$ tetap menjadi eksperimen portabilitas VLA terpisah; metode \textit{flow-matching} kontinu tidak disebut GRPO tanpa estimator probabilitas kebijakan yang tervalidasi. % [SRC-B2-11]

\subsection{Penambatan Visual, Penalaran Spasial, dan Instruksi Multilangkah}
\label{subsec:grounding_spasial}
Penambatan visual adalah pemetaan entitas yang disebut dalam instruksi ke objek atau wilayah pada pengamatan. Proses ini mencakup identifikasi kelas, atribut, dan instans yang benar di antara objek pengalih. Dalam \textit{vision-language navigation}, landmark serta relasi spasial pada instruksi perlu ditambatkan ke modalitas visual agar agen dapat memilih arah yang sesuai \citep{zhang2022explicit}. DroneVL Arena karena itu mengevaluasi dua keluaran yang berbeda: kecocokan instans target yang dideklarasikan terhadap \textit{ground truth}, serta terpenuhinya predikat spasial oleh pose UAV terhadap target. Setiap episode menyatakan kerangka acuan, toleransi geometris, dan waktu evaluasi bagi relasi seperti kiri, kanan, atau depan. % [SRC-B2-04]

Instruksi multilangkah menambahkan kendala urutan pada masalah penambatan. Pencapaian wilayah akhir tidak selalu berarti misi berhasil apabila subtujuan sebelumnya belum dipenuhi. Agen perlu membedakan instruksi yang telah selesai, instruksi yang menjadi fokus berikutnya, arah gerak yang relevan, dan kemajuan menuju sasaran. Kebutuhan tersebut menjadi dasar pendekatan \textit{self-monitoring} pada \textit{vision-language navigation} \citep{ma2019selfmonitoring}. DroneVL menggunakan riwayat tindakan yang terbatas dan ringkasan kemajuan misi sebagai bagian pengamatan kanonis agar keputusan berikutnya dapat mempertimbangkan konteks tanpa membuat masukan model bertumbuh tanpa batas. % [SRC-B2-05]

\subsection{\textit{Embodied Closed-Loop Agent}}
\label{subsec:closed_loop}
\textit{Embodied agent} menjalankan lingkar persepsi--keputusan--aksi--pengamatan. Berbeda dengan evaluasi gambar statis, tindakan pada navigasi \textit{closed-loop} mengubah keadaan dan pengamatan pada langkah berikutnya. Rotasi dapat membuka pandangan baru ke sasaran, sedangkan gerakan maju yang salah dapat mengurangi jarak bebas, menyebabkan tabrakan, atau menghabiskan anggaran misi. Karena itu, evaluasi harus mengukur konsekuensi rangkaian keputusan, bukan hanya ketepatan satu jawaban terhadap satu gambar.

DroneVL menerapkan horizon aksi pendek dan keputusan berulang. Setelah satu aksi atau satu durasi terbatas, sistem memperoleh pengamatan baru, membangun kembali konteks misi, dan meminta keputusan berikutnya dari \textit{backend}. Batas kecepatan dan durasi membatasi paparan serta perpindahan akibat satu keputusan, tetapi tidak menjamin bahwa akibatnya dapat dipulihkan; karena itu gerbang keselamatan tetap memeriksa jarak bebas, geofence, usia pengamatan, dan kondisi penghentian. Pengelola misi menentukan keadaan misi dan kondisi terminal, sedangkan pengendali tingkat rendah tetap sama pada seluruh \textit{backend}. Kesamaan ini mengendalikan perubahan pada jalur hilir, sementara perbedaan modalitas, resolusi, latensi, dan kapabilitas \textit{backend} dicatat sebagai faktor eksperimen, bukan diasumsikan setara.

\subsection{CoSyS-AirSim dan ROS~2}
\label{subsec:cosys_ros}
AirSim menyediakan simulasi fisik dan visual untuk kendaraan otonom \citep{shah2017airsim}, sedangkan CoSyS-AirSim memperluas kerangka tersebut untuk kebutuhan simulasi industri yang lebih kompleks \citep{jansen2023cosysairsim}. CoSyS-AirSim digunakan penelitian ini untuk membangun episode UAV multirotor dengan citra, kedalaman, segmentasi instans, dan lapisan anotasi. Pada kondisi utama, hanya citra RGB depan, kedalaman terdaftar, keadaan UAV, ringkasan kemajuan misi, dan instruksi yang memasuki \texttt{CanonicalObservation}. Segmentasi instans, label objek, dan transformasi objek merupakan data oracle yang hanya tersedia bagi pembangun manifest, \textit{baseline} oracle, dan evaluator; adaptor, pengelola misi, gerbang keselamatan, serta pengendali tidak dapat membacanya untuk memilih aksi. % [SRC-B2-12]

ROS~2 menyediakan middleware robotika yang berpusat pada simpul, topik, layanan, aksi, parameter, dan transformasi \citep{macenski2022ros2}. Jembatan penelitian menerbitkan citra kamera, \textit{CameraInfo}, odometri, IMU, GPS, data jarak, dan transformasi objek; jembatan juga menerima perintah kecepatan serta tujuan posisi lokal. Perbedaan kerangka koordinat antara simulator, ROS~2, dan badan UAV harus dikelola pada pembangun pengamatan serta adaptor pengendali. Adaptor VLM tidak boleh melakukan konversi koordinat secara mandiri karena konversi ganda dapat menghasilkan tindakan yang tidak konsisten dengan keadaan kendaraan. % [SRC-B2-06]

\subsection{Adaptor Agnostik-Model dan Aksi Terstruktur}
\label{subsec:adapter_aksi}
Adaptor agnostik-model menyembunyikan perbedaan permintaan dan respons setiap \textit{backend} di balik antarmuka yang tetap. Pada DroneVL, adaptor mendeklarasikan kapabilitas, menyiapkan permintaan dari \texttt{CanonicalObservation}, menjalankan inferensi, lalu mengembalikan respons mentah beserta metadata. Pengurai memiliki tanggung jawab tunggal untuk mendeserialisasi respons khusus-\textit{backend} menjadi objek aksi kanonis; kegagalan pada tahap ini dicatat sebagai kegagalan penguraian, bukan kegagalan adaptor. Adaptor tidak berlangganan topik ROS secara langsung, tidak mengendalikan UAV, tidak menentukan keberhasilan misi, dan tidak dapat melewati gerbang keselamatan. Batas tanggung jawab ini membuat penggantian \textit{backend} dapat dilakukan tanpa mengubah prosedur pengambilan sensor, pengendalian, maupun evaluasi.

Kontrak DroneVL membedakan perintah kendali dari usulan peristiwa misi. Perintah kendali yang dapat diteruskan menuju pengendali tingkat rendah adalah \texttt{MOVE\_RELATIVE}, untuk gerakan relatif berbatas; \texttt{ROTATE\_YAW}, untuk perubahan orientasi yaw; \texttt{HOVER\_AND\_OBSERVE}, untuk memperoleh pengamatan baru; dan \texttt{ABORT}, untuk penghentian terkendali. \texttt{DECLARE\_TARGET} serta \texttt{ADVANCE\_SUBGOAL} adalah usulan semantik kepada pengelola misi, bukan perintah penerbangan. Pengelola misi memverifikasi usulan tersebut terhadap predikat episode dan \textit{ground truth} evaluator sebelum mengubah keadaan misi.

Pengurai memeriksa struktur respons, bidang wajib, tipe data, nilai numerik berhingga, rentang tindakan, dan legalitas aksi pada keadaan misi saat ini. Satu \textit{prompt} perbaikan dapat digunakan apabila respons tidak valid; setelah $N_{\mathrm{parse}}$ kegagalan penguraian berturut-turut yang nilainya dibekukan pada manifest episode, pengelola misi menghentikan episode dengan alasan terminal \texttt{PARSE\_FAILURE}. Penghitung direset hanya setelah aksi kanonis yang valid diterima. Teks penjelasan tetap dicatat sebagai bukti, tetapi tidak digunakan sebagai parameter penerbangan.

\subsection{Batasan Keselamatan dan \textit{Benchmark} yang Dapat Direproduksi}
\label{subsec:keselamatan_benchmark}
Gerbang keselamatan pada DroneVL adalah mekanisme operasional untuk simulasi, bukan sertifikasi keselamatan penerbangan nyata. Gerbang tersebut memeriksa validitas skema, batas gerak dan yaw, geofence, ketinggian, usia data, jarak bebas dari informasi kedalaman atau rintangan, serta anggaran waktu dan keputusan. Setiap aksi dapat diterima, dipotong, diganti, atau ditolak dengan alasan yang dicatat. Log menyimpan aksi kanonis yang diusulkan, keputusan gerbang, aksi aktual, dan alasan intervensi agar keselamatan kebijakan mentah dapat dibedakan dari keselamatan sistem yang dibantu.

\textit{Benchmark} \textit{closed-loop} yang dapat direproduksi harus mendefinisikan lebih dari kumpulan gambar. Setiap episode memuat adegan, \textit{seed}, pose awal, instruksi, target, kondisi lingkungan, anggaran waktu dan keputusan, batasan, kondisi terminal, dan \textit{ground truth}. Manifest berversi juga memuat label kemampuan yang dapat bertumpang tindih, pemisahan pengembangan dan pengujian, serta versi kode, aset, model, \textit{prompt}, perangkat keras, dan log mentah. Split pengujian dibekukan sebelum penyetelan LoRA/GRPO dan dipisahkan dari pengembangan berdasarkan tata ruang, instans objek, komposisi instruksi, serta kondisi visual. Kebutuhan evaluasi kecerdasan spasial yang mempertimbangkan karakteristik navigasi UAV ditekankan oleh \textit{benchmark} VLM untuk UAV \citep{zhang2025skyready}. DroneVL Arena menghitung metrik dari log evaluator, bukan dari tabel yang diedit secara manual. % [SRC-B2-07]

%-----------------------------------------------------------------------------%
\section{Penelitian Terkait}
\label{sec:penelitian_terkait}
%-----------------------------------------------------------------------------%

\subsection{\textit{Vision-Language Model}, \textit{Embodied Agent}, dan Navigasi UAV}
\textit{Vision-language model} menunjukkan bahwa representasi visual dan bahasa dapat digunakan untuk memahami objek dan instruksi, tetapi representasi tersebut belum menentukan sendiri batas tindakan yang aman bagi robot \citep{radford2021clip,vinyals2015show}. Pada konteks UAV, van der Wal mengeksplorasi penggunaan VLM untuk eksplorasi aktif pada inspeksi anomali, yang menunjukkan relevansi penalaran visual--bahasa bagi tugas pengamatan dari udara \citep{wal2026active}. DroneVL mengambil posisi berbeda dengan menempatkan VLM di balik kontrak aksi yang terbatas, pengurai, dan gerbang keselamatan. Dengan demikian, evaluasi tidak hanya menilai apakah model menghasilkan respons semantik yang masuk akal, tetapi juga apakah respons tersebut dapat dinormalisasi dan dieksekusi dalam lingkar kendali yang tetap. % [SRC-B2-08]

Penelitian \textit{vision-language navigation} menekankan pelaksanaan instruksi dari pengamatan berurutan, penambatan visual, dan pemantauan kemajuan \citep{zhang2022explicit,ma2019selfmonitoring}. Tinjauan navigasi tujuan-objek pada UAV juga menunjukkan ketergantungan antara persepsi, pelokalan, perencanaan, dan kendali serta kebutuhan standardisasi kerangka evaluasi \citep{ayala2024uav}. \textit{Benchmark} \textit{Is your VLM Sky-Ready?} secara khusus mengarahkan evaluasi pada kecerdasan spasial untuk navigasi UAV \citep{zhang2025skyready}. DroneVL Arena sejalan dengan arah tersebut melalui tugas penambatan, penalaran spasial, pencarian, navigasi, dan instruksi multilangkah; pembeda utamanya adalah penggunaan satu kontrak pengamatan, aksi, pengendali, dan prosedur pencatatan untuk membandingkan \textit{backend} yang berbeda. % [SRC-B2-09]

\subsection{VLA Khusus UAV: CognitiveDrone}
\label{subsec:cognitivedrone_terkait}
CognitiveDrone merupakan penelitian yang secara langsung menerapkan \textit{vision-language-action} pada UAV. Lykov dkk. memperkenalkan model yang dilatih pada lebih dari 8.000 trajektori penerbangan simulasi untuk tugas pengenalan manusia, pemahaman simbol, dan penalaran; model menerima masukan visual orang-pertama serta instruksi teks, kemudian menghasilkan perintah aksi empat-dimensi secara waktu nyata \citep{lykov2025cognitivedrone}. Penelitian tersebut juga memperkenalkan CognitiveDroneBench sebagai \textit{benchmark} untuk tugas kognitif UAV dan varian CognitiveDrone-R1 yang menambahkan modul penalaran VLM sebelum kendali berfrekuensi tinggi.

Relevansi CognitiveDrone bagi penelitian ini terletak pada dua aspek. Pertama, ia menyediakan titik pembanding VLA yang asumsi perwujudannya sudah dekat dengan UAV. Kedua, ia menunjukkan bahwa evaluasi tindakan kognitif UAV memerlukan lebih dari pengenalan visual statis. Namun, hasil CognitiveDroneBench tidak diperlakukan sebagai hasil yang dapat dibandingkan langsung dengan DroneVL Arena karena lingkungan, episode, skema aksi, dan prosedur evaluasinya berbeda. Dalam DroneVL, CognitiveDrone ditempatkan sebagai satu kondisi pada eksperimen portabilitas VLA; keluarannya tetap harus menyatakan skema aksi, satuan, kerangka koordinat, serta horizon sebelum dapat dinormalisasi menjadi aksi kanonis.

Dengan posisi tersebut, DroneVL melengkapi, bukan menggantikan, CognitiveDrone. CognitiveDrone menilai kebijakan VLA khusus UAV dalam \textit{benchmark}-nya sendiri, sedangkan DroneVL menguji apakah kebijakan khusus UAV tersebut, OpenVLA-7B yang diadaptasi untuk UAV, dan OpenPI $\pi_{0.5}$ yang diadaptasi untuk UAV dapat melewati kontrak pengamatan, penguraian, gerbang keselamatan, dan evaluator yang sama. Pemisahan ini memungkinkan kegagalan karena format keluaran atau asumsi perwujudan dibedakan dari kegagalan navigasi pada setiap model.

\subsection{Celah Integrasi, Pengaman, dan \textit{Benchmark}}
CoSyS-AirSim dan ROS~2 menyediakan dua fondasi yang saling melengkapi, tetapi berada pada lapisan yang berbeda. CoSyS-AirSim memperluas AirSim melalui sensor, jenis kendaraan, dan lingkungan yang dapat diubah untuk simulasi industri serta pengujian algoritme navigasi \citep{jansen2023cosysairsim}. ROS~2 menyediakan arsitektur middleware berbasis simpul, topik, layanan, aksi, parameter, dan transformasi untuk menghubungkan komponen robotika \citep{macenski2022ros2}. Kedua karya tersebut menjadi dasar yang kuat untuk sensor, simulasi, dan komunikasi pada penelitian ini, tetapi kontribusi yang dilaporkan berfokus pada infrastruktur, bukan pada kontrak adaptor yang menormalkan keluaran VLM dan VLA heterogen menjadi aksi navigasi UAV yang dapat diaudit.

Pada lapisan model dan evaluasi, penelitian terkait juga menutup bagian masalah yang berbeda. CognitiveDrone menyediakan VLA khusus UAV beserta CognitiveDroneBench, dengan keluaran aksi empat-dimensi dari pengamatan orang-pertama dan instruksi \citep{lykov2025cognitivedrone}. Sementara itu, SpatialSky-Bench dalam \textit{Is your VLM Sky-Ready?} mengevaluasi kecerdasan spasial VLM untuk navigasi UAV melalui kategori persepsi lingkungan dan pemahaman adegan yang dibagi menjadi 13 subkategori \citep{zhang2025skyready}. Kedua karya tersebut menunjukkan bahwa VLA khusus UAV dan evaluasi kemampuan spasial VLM sama-sama diperlukan, tetapi memakai model, tugas, representasi keluaran, dan prosedur evaluasi yang dirancang untuk tujuan masing-masing.

Berdasarkan sintesis tersebut, celah yang ditangani DroneVL adalah lapisan integrasi dan evaluasi yang memungkinkan beberapa VLM serta VLA diuji pada kontrak pengamatan, aksi kanonis, pengendali tingkat rendah, batas keselamatan, episode, dan pencatatan yang sama. Pernyataan ini merupakan posisi penelitian yang diturunkan dari cakupan karya-karya di atas, bukan klaim bahwa karya terdahulu tidak memiliki nilai integrasi atau evaluasi. Celah tersebut diuji dengan melaporkan setiap tahap, yaitu inferensi, penguraian, keputusan gerbang keselamatan, aksi aktual, dan hasil misi, sehingga perbedaan kegagalan antarmodel dapat ditelusuri, bukan hanya dibandingkan melalui satu angka keberhasilan.

%-----------------------------------------------------------------------------%
\section{Posisi Penelitian}
\label{sec:posisi_penelitian}
%-----------------------------------------------------------------------------%
Tabel~\ref{tab:posisi_penelitian} merangkum posisi DroneVL terhadap kelompok penelitian terkait. Kontribusi penelitian bukan VLM baru atau klaim kelayakan penerbangan nyata, melainkan arsitektur integrasi dan evaluasi yang membuat perbandingan \textit{backend} dapat ditelusuri pada lingkungan UAV simulasi.

\begin{table}[h!]
    \centering
    \captionsetup{justification=justified, singlelinecheck=false}
    \caption{Posisi DroneVL terhadap kelompok penelitian terkait.}
    \label{tab:posisi_penelitian}
    \small
    \begin{tabularx}{\linewidth}{|>{\raggedright\arraybackslash}p{3.0cm}|>{\raggedright\arraybackslash}X|>{\raggedright\arraybackslash}X|>{\raggedright\arraybackslash}X|}
        \hline
        \textbf{Aspek} & \textbf{VLM dan navigasi berwujud} & \textbf{CoSyS-AirSim / ROS~2} & \textbf{DroneVL} \\
        \hline
        Pemahaman visual--bahasa & Tersedia dan bergantung pada model. & Hanya fondasi sensor. & Tersedia melalui \textit{backend} VLM. \\
        \hline
        Domain utama & Gambar statis atau agen darat pada banyak penelitian. & Simulasi kendaraan dan robot. & UAV multirotor simulasi. \\
        \hline
        Navigasi UAV \textit{closed-loop} & Tidak selalu menjadi fokus. & Menyediakan fondasi kendali. & Menjadi konfigurasi evaluasi utama. \\
        \hline
        \textit{Backend} VLM yang dapat diganti & Umumnya tidak dibahas. & Tidak disediakan. & Disediakan melalui adaptor. \\
        \hline
        Kontrak aksi dan gerbang keselamatan & Bergantung pada implementasi. & Kendali dasar tersedia. & Ditetapkan bersama dan dicatat. \\
        \hline
        Episode semantik dan evaluator berversi & Bergantung pada \textit{benchmark}. & Tidak disediakan sebagai arena semantik. & Disediakan oleh DroneVL Arena. \\
        \hline
    \end{tabularx}
\end{table}

Secara ringkas, DroneVL memiliki tiga kontribusi yang saling terkait. Pertama, penelitian ini merancang arsitektur agnostik-model yang memisahkan simulator, adaptor, pengurai, gerbang keselamatan, dan pengendali tingkat rendah. Kedua, DroneVL Arena mendefinisikan episode berversi, anotasi, pemisahan data, log, dan evaluator untuk tugas navigasi semantik UAV. Ketiga, penelitian ini menyiapkan perbandingan empiris yang melaporkan kualitas, keselamatan, efisiensi, aksi tidak valid, penggunaan sumber daya, dan latensi. Batas kontribusi tersebut adalah lingkungan simulasi dan kondisi eksperimen yang ditetapkan; penelitian ini tidak mensertifikasi keselamatan penerbangan nyata.
